\documentclass[a4paper, 11pt]{article}

% --- パッケージ ---
\usepackage[T1]{fontenc}
\usepackage[utf8]{inputenc}
\usepackage{geometry}
\usepackage{booktabs}
\usepackage{graphicx}
\usepackage{parskip}
\usepackage{titling}
\usepackage{xcolor}

% --- 余白を狭めに ---
\geometry{top=15mm, bottom=15mm, left=20mm, right=20mm}

% --- タイトルの上の余白を詰める ---
\setlength{\droptitle}{-2em}

% ==============================
\begin{document}

% --- ヘッダ部分 ---
\begin{center}
  \hrule height 1.5pt \vspace{6pt}
  {\LARGE\bfseries 作業報告書}\\[3pt]
  {\small 山田 太郎 / ○○部門 / \today}
  \vspace{6pt}\hrule height 1.5pt
\end{center}

% --- 本文 ---
\section{背景・目的}
ここに背景と目的を簡潔に書きます。

\section{内容・結果}
調査・作業の内容と結果を記載します。

% 表が必要な場合
\begin{center}
\begin{tabular}{@{}lcc@{}}
  \toprule
  項目 & 値A & 値B \\
  \midrule
  項目1 & 100 & 200 \\
  項目2 & 150 & 180 \\
  \bottomrule
\end{tabular}
\end{center}

\section{まとめ・所見}
結論や今後の対応などを書きます。

\end{document}
